\SeriesFrontHeading{The two volumes}
\noindent\textsc{Realization Calculus}\par\medskip
These volumes study realization under an explicit contract: which data are available, which transformations are allowed, and which objects count as answers. A compatible description, an admissible realization, and a recoverable source are different mathematical achievements. Their relation has to be proved.

\emph{Volume I, Komorebi Space}, constructs spaces of coherent families from declared systems of views. It develops extension obstructions, symmetry, descent, transport, and quantitative repair. Its analytic layers distinguish exact compatibility from membership in an intended normed realization, and a small residual from a small correction. Each conclusion retains the hypotheses that make it valid.

\emph{Volume II, Realization Boundaries}, examines the information that observations and completion can preserve or discard. A concrete family of countable abelian groups has one common profinite completion and one common unmarked finite-quotient system, while its abstract group types and effective presentations retain an infinite parameter. Subgroup classification, homomorphism rigidity, root structure, computability, and grammar escape measure that difference precisely.

The point of contact is the semantic comparison in Volume I: defining a category of coherent data does not establish that a functor from intended objects is faithful, full, or essentially surjective. Volume II supplies explicit categorical and computational tests in one group theoretic setting. Its examples preserve the universal property of the ambient limit. They show that the limit alone need not recover the chosen source image.

The volumes can be read independently. Read them together to move from the construction of coherent answers to the classification of what those answers determine. A finite resolution level in Volume I need not have a finite value set. The finite quotients of Volume II are finite groups. Likewise, escape from a fixed countable grammar is a relative theorem about that grammar and its semantics, rather than an assertion that the groups escape every possible system of views.

\vfill\noindent\textit{Komorebi P.}
\clearpage
\SeriesFrontHeading{A route through the pair}
\noindent The following connections give a short route across the two books. Section and chapter numbers within each volume retain their original mathematical numbering.

\begingroup\small\setlength{\tabcolsep}{5pt}\renewcommand{\arraystretch}{1.25}
\noindent\begin{tabularx}{\linewidth}{@{}p{.20\linewidth}XX@{}}
\toprule
\textbf{Question}&\textbf{Volume I: Komorebi Space}&\textbf{Volume II: Realization Boundaries}\\
\midrule
What is coherent?&Views and coherent realizations. Compatible evaluations determine a family.&Chapters 2 and 9. The prime-coordinate product and the intended dense source.\\
What do observations retain?&Finite observations and realized transformations. The comparison data specify what survives.&Chapters 3 and 9. Membership observations retain a spectrum; unmarked finite quotients forget the selector.\\
Which ambient points are intended?&Coherent points and Hilbert realizations. Compatibility and range membership are separate.&Chapters 4 and 9. Pure intermediate groups share an ambient product without becoming the same source.\\
Which maps lift?&Internal realization and intended objects. Faithfulness, fullness, and essential surjectivity are separate questions.&Chapter 9. Faithful completion maps, explicit non-lifts, and qualified reflection results.\\
What can a procedure supply?&Certificates and implementation boundaries. An existence proof and an executable repair need separate evidence.&Chapters 8 and 10. Oracle presentation semantics and escape from a fixed countable grammar.\\
\bottomrule
\end{tabularx}\endgroup

\section*{Formal companions}
Volume I carries its accepted Komorebi Space development: 162 production modules and four entrypoints, with formal counterparts for all sixteen principal theorem groups. This does not assign a separate Lean theorem to every illustration or printed proof route. Volume II carries its accepted Realization Boundaries development: 66 production modules and all forty-five registered contracts complete. Each volume's formal appendix explains its own scope.

This paired edition adds common presentation and reading connections. It preserves the mathematical bodies and the separately accepted proof sources. Its editorial checks establish source identity and publication quality. They do not constitute a new joint Lean replay. The source package records the accepted receipts, dependency trust boundaries, and the precise editorial changes.
\clearpage
